% NOT included in main.tex yet. After the user test, paste this section back
% into chapters/results.tex, just before \section{Threat Model}.

\section{User Study}
\label{sec:results-usability}

To find out whether people can actually use ChainBoard and understand its audit check, I ran a small usability study.

\paragraph{Participants and procedure.} \textit{tbd} people took part (\textit{tbd} students and \textit{tbd} developers; \textit{tbd} had used Trello or Jira before, \textit{tbd} had used a blockchain application). Each session took about 25 minutes on my laptop. Participants followed a printed sheet of 12 tasks in a prepared project (\Cref{tab:usability-tasks}) and were asked to think aloud. I did not help unless someone was stuck for more than two minutes, which I counted as a failure. Afterwards they filled in the System Usability Scale (SUS)~\cite{brooke1996sus} and two extra statements about the audit trail, on the same 1--5 scale, plus three open questions.

\begin{table}[tb]
    \centering
    \small
    \begin{tabular}{p{0.5cm}p{6.6cm}rr}
        \toprule
        \# & Task & Success & Median time (s) \\
        \midrule
        T1 & Sign in & \textit{tbd}/\textit{n} & \textit{tbd} \\
        T2 & Open the project and its board & \textit{tbd}/\textit{n} & \textit{tbd} \\
        T3 & Find the overdue task & \textit{tbd}/\textit{n} & \textit{tbd} \\
        T4 & Assign it to yourself & \textit{tbd}/\textit{n} & \textit{tbd} \\
        T5 & Move it to In progress & \textit{tbd}/\textit{n} & \textit{tbd} \\
        T6 & Add a comment & \textit{tbd}/\textit{n} & \textit{tbd} \\
        T7 & Create a task with priority and due date & \textit{tbd}/\textit{n} & \textit{tbd} \\
        T8 & Attach a file & \textit{tbd}/\textit{n} & \textit{tbd} \\
        T9 & Mark the task Done & \textit{tbd}/\textit{n} & \textit{tbd} \\
        T10 & Find who completed another task, and when & \textit{tbd}/\textit{n} & \textit{tbd} \\
        T11 & Verify the audit trail and explain the result & \textit{tbd}/\textit{n} & \textit{tbd} \\
        T12 & Try to change another person's task, explain why it fails & \textit{tbd}/\textit{n} & \textit{tbd} \\
        \bottomrule
    \end{tabular}
    \caption{Tasks of the user study, success rate and median completion time.}
    \label{tab:usability-tasks}
\end{table}

\paragraph{Task success.} \ag{Two or three sentences: overall success rate, which tasks were hardest and why (use what people said while thinking aloud).} For T11 and T12 a task counted as successful only if the participant could also explain the result in their own words, for example that the history is stored on a ledger and cannot be changed, or that only the assignee may change a task.

\paragraph{SUS score.} Each SUS answer was converted to a 0--100 score in the standard way: for odd items subtract 1 from the answer, for even items subtract the answer from 5, add the ten values and multiply by 2.5. \Cref{tab:sus} gives the results. An average SUS score is around 68~\cite{lewis2018sus}, so \ag{compare the mean with 68 in one sentence}.

\begin{table}[tb]
    \centering
    \small
    \begin{tabular}{lrrrr}
        \toprule
        Measure & Mean & SD & Min & Max \\
        \midrule
        SUS score (0--100) & \textit{tbd} & \textit{tbd} & \textit{tbd} & \textit{tbd} \\
        ``I understood what the result of Verify audit trail means'' (1--5) & \textit{tbd} & \textit{tbd} & \textit{tbd} & \textit{tbd} \\
        ``Being able to check a task's history makes me trust it more'' (1--5) & \textit{tbd} & \textit{tbd} & \textit{tbd} & \textit{tbd} \\
        \bottomrule
    \end{tabular}
    \caption{Questionnaire results ($n = \textit{tbd}$).}
    \label{tab:sus}
\end{table}

\paragraph{What participants said.} \ag{Summarize the open answers and think-aloud notes in a short paragraph or a small table: the two or three most common problems, what people liked, and one or two short quotes.}

\paragraph*{Changes made after the first session.} The study was used to improve the prototype while it was running. In the first session, the participant did not find how to set the priority and the due date of a new task (T7): the create-task bar showed only a title field and an ``Add task'' button, and the other fields were hidden behind a small ``More options'' link. After this session the create-task card was redesigned. The board now shows a single ``Add new task'' button, which opens a card that has a heading, a labelled title field, the priority as three visible buttons (low, medium, high), a visible due-date field, an optional description, and a ``Create task'' button with a hint when the title is missing; the task identifier is always generated automatically and is no longer shown. The result of the audit verification was also rewritten in plain language. The button ``Verify audit trail'' became ``Check that this history is genuine'', with one sentence explaining what it does; the result says, for example, ``This history is genuine: all 3 changes were found in the blockchain exactly as they are shown here'', the head digest is shown as a ``history code'' with a short explanation, and transaction identifiers moved under ``Technical details''. The following sessions used the new version. \ag{Say how many sessions used the old and the new version, and whether T7 became easier.}

\paragraph*{Changes made after the second session.} \ag{Start with what you observed in session 2 that led to these changes, e.g. where the participant got lost or hesitated.} After the second session the layout of the application was reorganized. A sidebar now lists the person's projects, so switching project takes one click from anywhere, and each project has a header with two tabs, ``Board'' and ``Members and history'', instead of separate pages and a ``Project settings'' button. A task no longer opens on its own page: it opens in a window over the board, with the status as three buttons under the title, a single line with priority, due date, assignee and task identifier, and below it the description, files, comments and activity. The separate ``Edit'' mode was removed; title, description, priority and due date are now changed in place and each change is saved as its own transaction. When a person tries to change something they are not allowed to, for example a task assigned to someone else, or to archive a task without being owner or admin, the interface now explains why instead of silently ignoring the click; the chaincode still enforces the same rules. Finally, the light theme was changed from grey to a warm off-white with near-black text and all text colours at or above the WCAG AA contrast of 4.5:1~\ag{cite WCAG 2.1}. \ag{Say which sessions used this version.}

With \textit{tbd} participants the study cannot give precise numbers, and most participants were students, not professional teams. It does show whether the main ideas, a familiar board and a visible audit check, make sense to people who did not build the system.
